%==========================================================
%  CHAPITRE 03 — INTRODUCTION GÉNÉRALE
%==========================================================
\chapter{Introduction générale}
\label{chap:intro}

\begin{flushright}
  \textit{\guillemotleft\ \todo{Citation inspirante en lien avec le projet}
  \guillemotright}\\
  \vspace{2pt}
  {\small\color{SecGrey}— \todo{Auteur de la citation}}
\end{flushright}

\vspace{1em}

%──────────────────────────────────────────────────────────
\section{Contexte général}
\label{sec:intro:contexte}
%──────────────────────────────────────────────────────────

\todo{Décrire en 2--3 paragraphes le contexte macro-économique,
sectoriel ou technologique dans lequel s'inscrit le projet.
Exemple : transformation numérique du secteur X, adoption croissante
du cloud, essor des solutions mobiles, etc.}

L'essor du numérique impose aux entreprises de tous secteurs une
transformation profonde de leurs systèmes d'information. Dans ce contexte,
\textbf{\CompanyName} \todo{décrire l'activité et les enjeux spécifiques}.

\vspace{0.5em}

\todo{Paragraphe 2 : préciser le périmètre du stage et l'équipe d'accueil.}

%──────────────────────────────────────────────────────────
\section{Problématique}
\label{sec:intro:problematique}
%──────────────────────────────────────────────────────────

\todo{Formuler clairement la problématique à laquelle répond le projet.
Une problématique est une question précise, pas un constat vague.
Exemple : « Comment automatiser le processus de validation des commandes
afin de réduire le délai de traitement de 48 h à moins de 2 h, tout en
garantissant la traçabilité des opérations ? »}

Face à ces enjeux, la question centrale à laquelle ce projet tente de
répondre est la suivante :

\begin{important}
  \centering
  \textit{« \todo{Formuler la problématique principale} »}
\end{important}

%──────────────────────────────────────────────────────────
\section{Objectifs du projet}
\label{sec:intro:objectifs}
%──────────────────────────────────────────────────────────

Ce projet de fin d'année vise à atteindre les objectifs suivants :

\begin{enumerate}
  \item \textbf{Objectif fonctionnel :}
        \todo{ex : développer une interface permettant la saisie et le
        suivi des commandes en temps réel.}

  \item \textbf{Objectif technique :}
        \todo{ex : concevoir une architecture REST scalable avec
        \tech{Spring Boot} et \tech{React}.}

  \item \textbf{Objectif qualité :}
        \todo{ex : atteindre une couverture de tests supérieure à 75\%
        et respecter les standards de codage de l'équipe.}

  \item \textbf{Objectif de déploiement :}
        \todo{ex : mettre en production sur l'infrastructure Docker
        de l'entreprise avant la fin du stage.}
\end{enumerate}

%──────────────────────────────────────────────────────────
\section{Méthodologie de travail}
\label{sec:intro:methodologie}
%──────────────────────────────────────────────────────────

\todo{Présenter brièvement la méthodologie adoptée (Agile/Scrum, cycle
en V, Kanban, etc.) et justifier ce choix au regard des contraintes du projet.}

Le projet a été conduit selon la méthodologie \textbf{\todo{Scrum / Agile / Cycle en V}},
en raison de \todo{justification : besoins évolutifs, petite équipe, délais courts, etc.}.
Les réunions hebdomadaires avec l'encadrant professionnel et les revues
de sprint bimensuelles ont permis d'assurer un suivi rigoureux.

%──────────────────────────────────────────────────────────
\section{Structure du rapport}
\label{sec:intro:structure}
%──────────────────────────────────────────────────────────

Ce rapport est organisé en \textbf{cinq chapitres principaux} :

\begin{description}[leftmargin=!, labelwidth=\widthof{\textbf{Chapitre 5}}]
  \item[\textbf{Chapitre 1}] \textit{Présentation de l'organisme d'accueil} —
    présente \CompanyName, son activité, son organigramme et
    l'environnement de travail.

  \item[\textbf{Chapitre 2}] \textit{Contexte général du projet} —
    analyse l'existant, détaille les besoins fonctionnels et
    non-fonctionnels et précise le cahier des charges.

  \item[\textbf{Chapitre 3}] \textit{Analyse et conception} —
    expose la méthodologie, les diagrammes UML et l'architecture
    technique retenue.

  \item[\textbf{Chapitre 4}] \textit{Réalisation} —
    détaille l'implémentation, les technologies utilisées et
    présente les principales interfaces.

  \item[\textbf{Chapitre 5}] \textit{Tests et résultats} —
    décrit la stratégie de tests, les résultats obtenus et
    les performances mesurées.
\end{description}

Une conclusion générale clôt le rapport en dressant le bilan
du projet et en proposant des perspectives d'évolution.

\clearpage
